\documentclass{article}

\usepackage[accepted]{icml2025}
\usepackage{microtype}
\usepackage{graphicx}
\usepackage{subcaption}
\usepackage{booktabs}
\usepackage{tabularx}
\usepackage{array}
\usepackage{amsmath}
\usepackage{amssymb}
\usepackage{url}
\usepackage{hyperref}

% Keep the ICML author/affiliation metadata, but suppress numeric affiliation
% markers after author names for this course report.
\makeatletter
\renewcommand{\@pa}[1]{%
  \ifcsname the@affil#1\endcsname
  \else
    \stepcounter{@affiliationcounter}%
    \newcounter{@affil#1}%
    \setcounter{@affil#1}{\value{@affiliationcounter}}%
  \fi}
\makeatother

\newcommand{\safeincludegraphics}[2][]{%
  \IfFileExists{#2}{\includegraphics[#1]{#2}}{%
    \fbox{\parbox[c][1.35in][c]{0.9\linewidth}{%
      \centering Missing figure: \texttt{#2}}}}}
\newcolumntype{Y}{>{\raggedright\arraybackslash}X}

\icmltitlerunning{Open-Vocabulary Detection and Grounding}

\begin{document}

\twocolumn[
\icmltitle{A Reproducible Comparison of Open-Vocabulary Object Detection\\
and Visual Grounding Models}

\begin{icmlauthorlist}
\icmlauthor{Zheng Renjie}{sustech}
\icmlauthor{Ling Hongyu}{sustech}
\icmlauthor{He Jiayang}{sustech}
\icmlauthor{Deng Haotian}{sustech}
\end{icmlauthorlist}

\icmlaffiliation{sustech}{Southern University of Science and Technology, Shenzhen, China}
\icmlcorrespondingauthor{Zheng Renjie}{12412840@mail.sustech.edu.cn}

\icmlkeywords{Open-Vocabulary Object Detection, Visual Grounding, OWL-ViT, Grounding DINO, YOLO-World}

\vskip 0.3in
]

% Force a compact running head; some XeTeX/Tectonic builds mis-measure the
% ICML template's title box and otherwise replace it with an error message.
\fancyhead[C]{\small\bf Open-Vocabulary Detection and Grounding}

\begin{abstract}
Open-vocabulary object detection replaces a detector's fixed label set with
arbitrary text prompts, while visual grounding localizes an object described by
a natural-language expression. We reproduce and compare three representative
pretrained systems: OWL-ViT Base, Grounding DINO Tiny, and YOLO-World v2 Small.
All models are evaluated without task-specific fine-tuning on the same
seeded random subset of 500 COCO val2017 images using the 80 category names as
prompts. Grounding DINO obtains the best COCO AP of 0.452, compared with 0.394
for YOLO-World and 0.270 for OWL-ViT. YOLO-World is substantially faster,
reaching 62.0 pipeline frames per second on an RTX 4060 Laptop GPU, versus 11.6
for OWL-ViT and 1.5 for Grounding DINO. We additionally evaluate genuine
referring-expression localization on a 100-expression RefCOCO validation
subset. Grounding DINO, OWL-ViT, and YOLO-World reach top-1
Accuracy@IoU$\geq0.5$ of 59\%, 56\%, and 46\%, respectively. Threshold and
non-maximum suppression ablations show that aggressive filtering improves
visual clarity but reduces recall and quantitative accuracy. The resulting
codebase provides reproducible data preparation, evaluation, timing, and
visualization scripts and exposes a clear accuracy--efficiency trade-off among
the three model families.
\end{abstract}

\section{Introduction}

Conventional object detectors predict from a category vocabulary fixed during
training. Adding a new concept generally requires collecting annotations and
retraining or fine-tuning the model. Open-vocabulary object detection (OVD)
instead conditions localization on text, enabling prompts such as ``remote
control,'' ``the red car,'' or concepts not explicitly represented by a fixed
output head. Visual grounding is closely related but more demanding: the input
is often a full referring expression whose attributes and spatial relations
must identify one instance among several candidates.

This project studies how readily available pretrained models behave under a
common, reproducible protocol. We compare OWL-ViT
\citep{minderer2022owlvit}, Grounding DINO \citep{liu2024groundingdino}, and
YOLO-World \citep{cheng2024yoloworld}. These systems span three useful design
points: a simple vision-transformer transfer approach, a tightly fused
language-conditioned transformer detector, and a real-time YOLO-style
detector.

Our work makes four practical contributions:
\begin{enumerate}
  \item a unified evaluation on the same fixed random 500-image COCO subset;
  \item synchronized runtime and peak GPU-memory measurements;
  \item a separate RefCOCO experiment that tests phrase-level grounding rather
  than only category-name detection; and
  \item threshold and NMS ablations that quantify common deployment choices.
\end{enumerate}
We do not claim a new detector. The goal is accurate reproduction, controlled
comparison, and analysis of when each approach works or fails.

\section{Related Work}

\paragraph{Vision-language pretraining.}
CLIP learns transferable visual representations by matching images and natural
language at scale \citep{radford2021clip}. Its shared embedding space makes text
an effective interface for downstream recognition, but image-level alignment
does not by itself provide object boxes.

\paragraph{Open-vocabulary detection.}
OWL-ViT transfers a contrastively pretrained vision transformer to
text-conditioned detection with minimal architectural changes
\citep{minderer2022owlvit}. It embeds each prompt and scores image tokens
against the resulting text features. Grounding DINO builds on transformer
detection ideas, including DETR \citep{carion2020detr}, and introduces
language-guided query selection plus cross-modal fusion
\citep{liu2024groundingdino}. YOLO-World targets real-time operation through a
re-parameterizable vision-language path aggregation network and a
prompt-then-detect workflow \citep{cheng2024yoloworld}.

\paragraph{Referring-expression comprehension.}
RefCOCO and related datasets pair COCO objects with expressions intended to
distinguish a target from other objects in the same image
\citep{yu2016referring}. Unlike category detection, success may require color,
position, attributes, or context. The conventional metric counts a prediction
as correct when its intersection-over-union (IoU) with the target box exceeds
0.5.

\section{Approach}

\subsection{Unified Text-Conditioned Pipeline}

Given an image $I$ and prompts $T=\{t_1,\ldots,t_K\}$, each model produces
candidate boxes, text labels, and confidence scores:
\begin{equation}
  f_\theta(I,T) = \{(b_i, \hat{t}_i, s_i)\}_{i=1}^{N}.
\end{equation}
We preserve each model's official preprocessing and pretrained weights. No
training or fine-tuning is performed. Boxes are converted to COCO's
$(x,y,w,h)$ format and evaluated with \texttt{pycocotools}.

\subsubsection*{OWL-ViT Base}
We use \texttt{google/owlvit-base-patch32}. For COCO, all 80 category names are
passed as independent text queries. Predictions below 0.01 are removed and the
top 100 are retained. Extra NMS is disabled in the main experiment because the
model's dense hypotheses are already ranked and our ablation showed that
additional suppression reduced AP.

\subsubsection*{Grounding DINO Tiny}
We use \texttt{IDEA-Research/grounding-dino-tiny}. COCO class names are joined
into a period-separated caption. Both box and text thresholds are 0.20, and
the top 100 matched predictions are retained. The model is evaluated in FP32
because this configuration was stable on the Windows test environment.

\subsubsection*{YOLO-World v2 Small}
We use \texttt{yolov8s-worldv2.pt} through the Ultralytics interface
\citep{jocher2023ultralytics}. The 80 COCO class embeddings are computed before
inference. Images are resized according to the model's 640-pixel input setting,
with confidence threshold 0.001, internal NMS IoU 0.7, and at most 100 boxes.

\subsection{Evaluation Metrics}

For COCO detection, we report mean average precision over IoU thresholds
$0.50{:}0.05{:}0.95$ (AP), AP50, AP75, scale-specific AP, and average recall
with at most 100 detections (AR100) \citep{lin2014coco}. IoU is
\begin{equation}
  \operatorname{IoU}(B,\hat B)=
  \frac{|B\cap\hat B|}{|B\cup\hat B|}.
\end{equation}

For RefCOCO, one expression is provided at a time and only the
highest-scoring box is evaluated. We report
\begin{equation}
  \operatorname{Acc@}\tau =
  \frac{1}{M}\sum_{j=1}^{M}
  \mathbb{1}\left[\operatorname{IoU}(B_j,\hat B_j)\geq\tau\right]
\end{equation}
for $\tau\in\{0.5,0.75\}$, together with mean IoU.

\subsection{Runtime Measurement}

Each model receives one warm-up image that is excluded from timing. CUDA is
synchronized around preprocessing, inference, and postprocessing for OWL-ViT
and Grounding DINO. YOLO-World's internal stage timings are accumulated.
File download and disk image loading are excluded. We report pipeline FPS and
PyTorch's peak allocated GPU memory. These numbers characterize our specific
hardware and software configuration rather than universal model speed.

\section{Experimental Setup}

\paragraph{Hardware and software.}
Experiments run on an NVIDIA GeForce RTX 4060 Laptop GPU with 8 GB VRAM.
The environment uses Python 3.10, PyTorch 2.11.0 with CUDA 12.8,
Torchvision 0.26.0, Transformers 5.5.3, Ultralytics 8.4.70, and
pycocotools 2.0.11.

\paragraph{COCO subset.}
We sample 500 images without replacement from COCO val2017 using Python's
fixed random seed 308. The sorted selected image IDs are stored with every
metrics file, so all models see exactly the same images. The 80 COCO categories
serve as text prompts. This is more representative than selecting the first
100 image IDs, but it remains a subset rather than a full COCO validation run.

\paragraph{RefCOCO subset.}
We use the first 100 region rows exposed by the
\texttt{lmms-lab/RefCOCO} validation mirror, with original COCO train2014
images and target boxes. Each row contains two to four human expressions. We
use the first expression per region to avoid repeatedly weighting the same
target. This experiment is a small reproducible diagnostic and should not be
compared directly with results on the complete official RefCOCO validation
split.

\section{Experimental Results}

\subsection{Open-Vocabulary Detection on COCO}

\begin{table*}[t]
\caption{Zero-shot detection on the fixed random 500-image COCO val2017
subset. Best values are bold.}
\label{tab:coco}
\centering
\small
\begin{tabular}{lrrrrrrrr}
\toprule
Model & AP & AP50 & AP75 & AP$_S$ & AP$_M$ & AP$_L$ & AR100 \\
\midrule
OWL-ViT Base & 0.270 & 0.428 & 0.287 & 0.134 & 0.303 & 0.454 & 0.487 \\
Grounding DINO Tiny & \textbf{0.452} & \textbf{0.589} & \textbf{0.498} &
\textbf{0.325} & \textbf{0.497} & 0.549 & 0.548 \\
YOLO-World v2 Small & 0.394 & 0.549 & 0.427 & 0.234 & 0.452 &
\textbf{0.556} & \textbf{0.561} \\
\bottomrule
\end{tabular}
\end{table*}

Table~\ref{tab:coco} and Figure~\ref{fig:coco} show a clear ranking in AP.
Grounding DINO exceeds YOLO-World by 5.8 AP points and OWL-ViT by 18.2 points.
Its advantage is largest for small objects: AP$_S$ is 0.325 versus 0.234 and
0.134. This supports the qualitative observation that small or cluttered
instances are a major weakness of OWL-ViT. YOLO-World has the highest AR100
and AP$_L$, indicating strong proposal coverage and localization for larger
objects, although its lower overall AP suggests more ranking or classification
errors than Grounding DINO.

\begin{figure}[t]
  \centering
  \safeincludegraphics[width=\linewidth]{figures/coco_metrics.pdf}
  \caption{COCO detection accuracy on the shared 500-image subset.}
  \label{fig:coco}
\end{figure}

\subsection{Accuracy--Efficiency Trade-off}

\begin{table}[t]
\caption{Measured efficiency on the same GPU. Pipeline time excludes disk I/O.}
\label{tab:efficiency}
\centering
\small
\begin{tabular}{lrrr}
\toprule
Model & ms/image & FPS & Peak MB \\
\midrule
OWL-ViT Base & 86.6 & 11.6 & 351 \\
Grounding DINO Tiny & 680.5 & 1.5 & 2356 \\
YOLO-World v2 Small & \textbf{16.1} & \textbf{62.0} & 716 \\
\bottomrule
\end{tabular}
\end{table}

YOLO-World is the practical speed winner (Table~\ref{tab:efficiency} and
Figure~\ref{fig:efficiency}). It is 5.4 times faster than OWL-ViT and over 42
times faster than Grounding DINO in pipeline FPS. Grounding DINO uses the most
memory and is the slowest, but provides the highest AP. OWL-ViT uses the least
allocated GPU memory, although its accuracy is lowest. Thus no model dominates:
Grounding DINO is preferable when accuracy is primary, YOLO-World when latency
is constrained, and OWL-ViT when a simple low-memory transformer baseline is
desired.

\begin{figure}[t]
  \centering
  \safeincludegraphics[width=\linewidth]{figures/efficiency.pdf}
  \caption{Pipeline throughput and peak allocated GPU memory.}
  \label{fig:efficiency}
\end{figure}

\subsection{Referring-Expression Grounding}

\begin{table}[t]
\caption{Top-1 localization on 100 RefCOCO validation expressions.}
\label{tab:refcoco}
\centering
\small
\resizebox{\columnwidth}{!}{%
\begin{tabular}{lrrr}
\toprule
Model & Acc@0.5 & Acc@0.75 & Mean IoU \\
\midrule
OWL-ViT Base & 0.56 & 0.50 & 0.513 \\
Grounding DINO Tiny & \textbf{0.59} & \textbf{0.56} & \textbf{0.591} \\
YOLO-World v2 Small & 0.46 & 0.41 & 0.468 \\
\bottomrule
\end{tabular}
}
\end{table}

Grounding DINO remains strongest on phrase-level localization
(Table~\ref{tab:refcoco}, Figure~\ref{fig:refcoco}), consistent with its
explicit cross-modal fusion and grounded pretraining. OWL-ViT is only three
points behind at IoU 0.5 and remains a competitive zero-shot grounding
baseline. YOLO-World is faster but loses 13 points relative to Grounding DINO,
suggesting that longer relational descriptions are harder for its
efficiency-oriented text-conditioned detector.

\begin{figure}[t]
  \centering
  \safeincludegraphics[width=\linewidth]{figures/refcoco_grounding.pdf}
  \caption{RefCOCO top-1 Accuracy@IoU thresholds.}
  \label{fig:refcoco}
\end{figure}

\subsection{Threshold Sensitivity and NMS}

The three model APIs expose differently calibrated scores, so a numerical
threshold of 0.2 does not have the same meaning across models. We therefore
vary each threshold only within a model, using saved predictions. As shown in
Figure~\ref{fig:threshold}, AP and AR decrease monotonically as detections are
removed. For OWL-ViT, increasing the threshold from 0.01 to 0.20 reduces AP
from 0.270 to 0.211 and AR100 from 0.487 to 0.265. Grounding DINO drops from
0.452 to 0.400 AP between 0.20 and 0.40. YOLO-World drops from 0.394 to 0.349
between 0.001 and 0.25. Higher thresholds can still be useful for uncluttered
visual demonstrations, but they are harmful when recall matters.

\begin{figure}[t]
  \centering
  \safeincludegraphics[width=\linewidth]{figures/threshold_sensitivity.pdf}
  \caption{Within-model score-threshold sensitivity on the COCO subset.}
  \label{fig:threshold}
\end{figure}

We also evaluated class-wise NMS for OWL-ViT on the earlier 100-image subset.
At NMS IoU 0.5, AP decreased from 0.336 to 0.323 and AR100 from 0.520 to
0.475. NMS visually removes duplicate boxes, but overlapping hypotheses can be
valid in crowded scenes. We therefore disable additional NMS for the main
OWL-ViT quantitative comparison and use it only for selected visualizations.

\subsection{Qualitative Analysis}

Figure~\ref{fig:qualitative} compares representative OWL-ViT and Grounding DINO
outputs. Both models handle common category nouns and multi-word prompts.
Grounding DINO generally produces tighter boxes and fewer missed small objects.
Common failure modes include duplicate low-confidence boxes, confusion between
nearby objects, failure to resolve spatial relations, and missed small
accessories. Prompt wording also matters: a broad noun can activate several
instances, while a detailed expression can either disambiguate the target or
reduce confidence if its phrasing differs from pretraining language.

\begin{figure*}[t]
  \centering
  \begin{subfigure}{0.24\textwidth}
    \safeincludegraphics[width=\linewidth]{figures/owlvit_workspace.jpg}
    \caption{OWL-ViT: workspace}
  \end{subfigure}
  \begin{subfigure}{0.24\textwidth}
    \safeincludegraphics[width=\linewidth]{figures/grounding_dino_workspace.jpg}
    \caption{Grounding DINO}
  \end{subfigure}
  \begin{subfigure}{0.24\textwidth}
    \safeincludegraphics[width=\linewidth]{figures/owlvit_traffic.jpg}
    \caption{OWL-ViT: traffic}
  \end{subfigure}
  \begin{subfigure}{0.24\textwidth}
    \safeincludegraphics[width=\linewidth]{figures/grounding_dino_traffic.jpg}
    \caption{Grounding DINO}
  \end{subfigure}
  \caption{Qualitative open-vocabulary detections. Display thresholds and NMS
  are chosen for readability and differ from the quantitative protocol.}
  \label{fig:qualitative}
\end{figure*}

\section{Limitations}

First, 500 COCO images reduce computation but produce less stable estimates
than the full 5,000-image validation set. Although fixed random sampling avoids
the obvious ordering bias of the original first-100 protocol, subset variance
remains. Second, COCO's 80 names are likely represented in the models'
large-scale pretraining data. The experiment therefore measures zero-shot
transfer without COCO-specific fine-tuning, not guaranteed generalization to
truly unseen concepts.

Third, the RefCOCO experiment uses only 100 mirrored validation rows and one
expression per target. A full evaluation should use every official split and
all expressions. Fourth, score calibration differs across architectures;
threshold sensitivity is meaningful within a model but cannot fully equalize
precision--recall operating points across models. Finally, runtime depends on
software versions, input resolution, precision, and hardware. In particular,
Grounding DINO was kept in FP32 for stability, while OWL-ViT used FP16 on GPU.

\section{Conclusion}

We reproduced three open-vocabulary detection systems and evaluated them under
a common protocol. Grounding DINO Tiny provides the best detection and
grounding accuracy, reaching 0.452 COCO AP and 59\% RefCOCO Acc@0.5.
YOLO-World v2 Small offers the strongest efficiency--accuracy compromise,
reaching 0.394 AP at 62.0 pipeline FPS. OWL-ViT is compact and simple but
trails on COCO, especially for small objects. The experiments also show that
postprocessing choices are not cosmetic: high confidence thresholds and extra
NMS can make figures cleaner while reducing AP and recall.

Future work should evaluate the full COCO and RefCOCO splits, add LVIS rare
categories or self-collected non-COCO objects to better test unseen concepts,
and compare prompt paraphrases systematically. Calibration or per-model
precision--recall operating-point selection would also improve deployment
comparisons.

\section*{Team Contributions}

All four members contributed equally (25\% each).

\bibliography{example_paper}
\bibliographystyle{icml2025}

\end{document}
